% Compact 4-page main report version. Compile in the official ACL template directory with acl.sty present.
\documentclass[11pt]{article}
\usepackage{acl}
\usepackage{times}
\usepackage{latexsym}
\usepackage[T1]{fontenc}
\usepackage[utf8]{inputenc}
\usepackage{microtype}
\usepackage{inconsolata}
\usepackage{booktabs}
\usepackage{amsmath}
\usepackage{graphicx}
\usepackage{array}
\usepackage{listings}
\usepackage{xcolor}
\usepackage{float}
\usepackage{placeins}
\usepackage{adjustbox}
\usepackage{url}
\lstset{basicstyle=\ttfamily\footnotesize, breaklines=true, breakatwhitespace=true, columns=fullflexible, keepspaces=true, frame=single, xleftmargin=0.5em, xrightmargin=0.5em}
\sloppy

\title{Do LLM Judges Prefer LLM-Generated Plans?\\
Order Effects and Criterion Decomposition in Pairwise LLM Evaluation}

\author{Cameron Caputa \\
}

\begin{document}
\maketitle

\begin{abstract}
LLMs are increasingly used as judges for open-ended NLP artifacts, but pairwise judgments may reflect position, source, or presentation cues rather than substantive quality. We test whether LLM judges prefer directly LLM-generated trail-running plans over programmatically structured controls when both are structurally matched and rendered through the same judge-facing schema. Across 250 matched pairs, four Qwen/Gemma judges, five runs, and AB/BA order swaps, LLM-source plans win only 36.53\% of 10,000 forced judgments. The dominant effect is positional: LLM-source plans win 57.54\% as Plan A but only 15.52\% as Plan B. Marker-level evaluation favors programmatic controls on every criterion in the order-balanced aggregate, but marker-by-order and marker-by-athlete-band splits show that several criteria are themselves context- and position-sensitive. Structural matching therefore improves source-preference studies but does not make pairwise LLM judgments reliable; order swaps, criterion-by-order analysis, athlete-band stratification, and checks for presentation leakage should be core design requirements.

Code is available at: \url{https://github.com/MauriceCC04/llm-judge-self-preference}
\end{abstract}

\section{Introduction}

LLM-as-judge evaluation is widely used for open-ended NLP outputs because many artifacts lack a single reference answer. However, pairwise LLM judging can reflect position bias, verbosity bias, self-enhancement bias, or source cues rather than substantive quality \citep{zheng2023llmasajudge,gu2024language}. This paper studies a narrower source-preference question: when artifacts are structurally matched and source-masked, do LLM judges prefer outputs produced by direct LLM-source generation over outputs produced by a programmatic structural generator?

We test this question in trail-running training-plan generation. Trail-running plans provide a useful testbed because they combine objectively describable structure---weekly load, rest days, hard-day spacing, and long-run placement---with subjective presentation qualities such as clarity, specificity, and rationale quality. The goal is not to determine which generator produces objectively better coaching plans. Instead, we ask whether a common LLM-as-judge setup can support a source-preference claim once structure, order, and presentation cues are controlled. Because the programmatic arm may include LLM-written narrative fields, the comparison is direct LLM-source generation versus programmatic structural generation, not fluent LLM prose versus empty text.

This study makes three contributions. First, it provides a controlled source-preference test using 250 structurally matched LLM-vs-programmatic pairs and 10,000 order-balanced pairwise judgments. Second, it shows that the main threat to interpretation is order sensitivity: LLM-source plans win 57.54\% as Plan A but only 15.52\% as Plan B. Third, it decomposes the result by athlete band and explicit evaluation markers, showing that the aggregate programmatic advantage is strongest for beginner profiles and that programmatic controls win every marker in aggregate, while several marker-level preferences change across athlete band and presentation order. Together, the results suggest that structural matching can make source-preference comparisons more controlled, but it cannot by itself make pairwise LLM judgments stable or reliable.

\section{Related Work}

LLM-as-judge evaluation uses models to score or compare open-ended artifacts when reference answers are unavailable or incomplete. MT-Bench and Chatbot Arena show that pairwise LLM judging can scale evaluation, but they also identify position bias, verbosity bias, self-enhancement bias, and reasoning limitations \citep{zheng2023llmasajudge}. Surveys of LLM-as-judge systems similarly emphasize prompt sensitivity, source cues, position effects, and aggregation choices \citep{gu2024language}. These studies motivate treating the judge itself as an experimental object rather than as a neutral measuring device.

Reliability work argues that repeated judgments are needed because single LLM judgments may not be stable across samples or settings \citep{schroeder2025trustllmjudgmentsreliability}. This is especially important for pairwise judging: AB/BA swaps can balance the aggregate source estimate while individual decisions remain unstable. We therefore distinguish an order-balanced aggregate preference from pair-level reliability across swapped presentations.

Criteria-decomposed evaluation methods such as HD-Eval motivate breaking holistic preferences into explicit dimensions rather than treating a single overall judgment as self-explanatory \citep{liu2024hdevalaligninglargelanguage}. Related work on learner-created artifacts and creative writing similarly shows that LLM evaluation quality can vary across rubric dimensions, domains, and prompt formats \citep{tian2024evaluating,kim2025llms}. Our marker pass adopts this diagnostic motivation for training plans: it tests whether the source effect appears in structure-sensitive criteria, such as recovery safety and load progression, or in presentation-sensitive criteria, such as clarity and explanation quality.

The gap for this paper is narrower than general LLM-as-judge benchmarking. Prior work establishes that LLM judges can be biased or unstable; we ask whether a source-preference claim still holds when the compared artifacts are matched on visible structure, presented through the same schema, evaluated in both AB and BA orders, and decomposed into explicit criteria. Crowd-comparison work instead tries to improve judges with richer anchors \citep{zhang2025crowd}; our goal is complementary, auditing a standard pairwise setup rather than proposing a new judge architecture.

\section{Experiments}

The final candidate pool contains 192 Qwen-source plans, 192 Gemma-source plans, and 640 programmatic plans. The programmatic arm samples structural constraints such as hard-day count, rest-day placement, recovery capacity, race phase, and load. A \emph{fixture} is one experimental cell: an athlete band crossed with readiness, recovery capability, and race phase. Athlete band is treated as a substantive matching variable because it changes both the generation context and what counts as an appropriate weekly structure. An early generation audit found substantial within-cell duplication, so the conditioning space was expanded with athlete profiles from beginner/low-volume to high-volume competitive runners.

Each LLM-source plan is matched to one programmatic plan using a source-neutral structural score over plan length, total minutes, rest days, hard days, quality days, long runs, duration plausibility, load-band fit, and hard-day spacing. The structural score is a matching aid, not a quality label: it measures visible weekly-plan similarity under fixture-conditioned expectations. Candidate pairs must share fixture ID, athlete band, readiness, recovery capability, race phase, plan length, and style, and must fall within a structural-score tolerance of 2.0. The matcher then selects a target-cardinality minimum-cost matching of 250 pairs using weighted structural feature distances and soft caliper penalties, where a soft caliper penalizes large feature gaps without forbidding them. The final matched set has mean structural-score gap 0.3076. For scale, a random same-fixture one-to-one baseline preserving the final per fixture counts has mean gap about 7.30, so the selected matching reduces structural-score gap by 95.8\% relative to random same-cell pairing.

Both plans are rendered through the same canonical masked judge-facing schema, which removes source/provenance metadata and preserves day-level structure plus short workout and purpose summaries. The frozen evaluation manifest stores source family, plan IDs, order, judge, run, and structural-score metadata outside the prompt, while the judge sees only neutral \texttt{plan\_a} and \texttt{plan\_b}. This is structural-score matching rather than full quality equalization: unmeasured coaching logic, wording style, and residual presentation cues may still differ between sources.

Four local judges, \texttt{qwen\_7b\_judge}, \texttt{qwen\_14b\_judge}, \texttt{gemma\_4b\_judge}, and \texttt{gemma\_12b\_judge}, evaluate each matched pair across five runs and two AB/BA orders, producing 10,000 forced judgments. The primary pairwise task uses \texttt{trailtraining.llm.soft\_eval.compare\_plans}: judges receive two anonymized canonical JSON plans labeled only as Plan A and Plan B and return a preferred plan, reasoning, and plan-specific advantages. The exact natural-language wrapper is inherited from the external \texttt{trailtraining} dependency rather than stored as a project-local prompt, so the reproducible evaluation object is the archived masked payload, judge configuration, and returned decision fields. A second pass evaluates nine \emph{markers}, or criterion-level preferences---plan coherence, training specificity, load progression, recovery safety, quality session design, endurance development, readiness alignment, clarity/actionability, and explanation quality---allowing ties and producing 90,000 marker decisions; the full marker prompt is reproduced in Appendix~\ref{app:marker-prompt}. Because judgments repeat the same matched pairs across judges, runs, and orders, row-level win rates are descriptive rather than independent-sample inference. We therefore treat matched-pair behavior, order-swapped consistency, marker decomposition, athlete-band stratification, and leakage-filtered sensitivity as the main robustness checks.

LLM-source plans used direct source generation at \(T=0.7\) followed by an explainer stage at \(T=0.0\). Programmatic plans had no LLM source-generation stage, but their LLM-assisted explanatory fields were generated at \(T=0.0\); all pairwise and marker judging also used \(T=0.0\), so repeated runs are not a judge-temperature sweep.

\section{Results}

Appendix Table~\ref{tab:summary} summarizes the main empirical findings. Overall, LLM-source plans do not win against structurally matched programmatic controls: judges select them in 36.53\% of forced judgments, while programmatic controls win 63.47\% (Appendix Table~\ref{tab:primary-app}). This aggregate hides substantial athlete-band heterogeneity (Appendix Tables~\ref{tab:athlete-app} and~\ref{tab:marker-band-app}). LLM-source plans win only 12.50\% for A1 beginner fixtures and 24.55\% for A2 developing fixtures, but reach 51.08\% for A3 advanced fixtures and 53.32\% for A4 competitive fixtures. Aggregating across markers, LLM-source win rates remain below 50\% within every athlete band, but the marker-by-band table shows that this aggregate masks criterion-specific reversals for advanced and competitive athletes: training specificity crosses parity for A3/A4 (58.72\% and 59.09\%), and quality-session design rises from 11.50\% for A1 to 51.06\% for A4. By contrast, readiness alignment and recovery safety remain programmatic-favoring across bands. Thus, the athlete-band effect reflects a shift in which criteria become competitive, not a uniform LLM-source advantage.

The dominant methodological result is the order effect (Appendix Tables~\ref{tab:order-main} and~\ref{tab:order-effect-app}). LLM-source plans win 57.54\% of judgments when shown as Plan A but only 15.52\% when shown as Plan B, a 42.02 percentage-point swing. AB/BA balancing prevents a trivial confound between source and position, but it does not make the judge's decision stable: many pair-level outcomes reverse when the same plans are swapped. Pair-run consistency shows the same problem (Appendix Tables~\ref{tab:diagnostics-main} and~\ref{tab:pair-run-app}). Here, source-consistent means the same source wins in both swapped orders, while position-consistent means the same screen position wins after the swap. Across matched pair--judge--run units, programmatic plans are source-consistent winners in 39.46\%, LLM-source plans are source-consistent winners in only 12.52\%, and Plan A is the position-consistent winner in 45.02\%. The experiment therefore estimates an order-balanced aggregate preference, not a stable pairwise decision mechanism.

Marker-level evaluation reinforces rather than reverses the aggregate result, but the stratified tables change how the marker results should be read (Appendix Tables~\ref{tab:diagnostics-main},~\ref{tab:markers-app},~\ref{tab:marker-order-app}, and~\ref{tab:marker-band-app}). Programmatic controls win every marker in the order-balanced aggregate. However, marker rates are not stable criterion-level preferences. Clarity/actionability, plan coherence, and recovery safety show extreme Plan-A sensitivity: LLM-source win rates fall from 85.40\%, 76.16\%, and 76.03\% when the LLM plan is Plan A to 1.21\%, 1.30\%, and 5.92\% when it is Plan B. Training specificity is the opposite caution: its aggregate near-parity rate of 48.54\% masks a reversal from 33.50\% as Plan A to 63.33\% as Plan B. Explanation quality remains programmatic-favoring in both orders, despite the programmatic arm's LLM-assisted explanatory fields. Across athlete bands, LLM-source plans become much more competitive in A3/A4 for training specificity, endurance development, and quality-session design, but not for readiness alignment or recovery safety. Thus, criterion decomposition is useful, but marker aggregates should be treated as order-balanced and band-averaged summaries rather than stable latent construct measurements. A post hoc leakage-filtered sensitivity analysis is stable: LLM-source plans win 34.44\% of clean-subset judgments, compared with 36.53\% overall (Appendix Table~\ref{tab:leakage-app}).


\section{Discussion and Limitations}

The main result is methodological rather than domain-specific. Structural matching and shared schemas reduce obvious confounds, but they do not make pairwise LLM judgments stable. AB/BA balancing protects the aggregate estimate from source-position imbalance, yet individual decisions remain highly position-sensitive. Judge-level order swings range from 20.00 to 66.00 percentage points, so the effect is not reducible to one idiosyncratic judge. This matters for NLP evaluation settings where small numbers of pairwise LLM judgments are used for model selection, leaderboard ranking, or artifact evaluation. A source-preference study that reports only one order per pair could reach a very different conclusion depending on which source appears first.

The athlete-band pattern also matters for interpretation. The aggregate programmatic advantage is driven mainly by beginner and developing profiles, where safe weekly structure may be more constrained: conservative total load, rest-day placement, and hard-day spacing leave less room for stylistic variation. The marker-by-band results clarify this mechanism: LLM-source plans gain most from A1 to A4 on training specificity, quality-session design, and endurance development, where richer contextual adaptation and more complex workouts may help. They do not gain similarly on recovery safety or readiness alignment, where explicit structural constraints remain important. This supports a more nuanced interpretation: direct LLM generation may help with flexible sport-specific tailoring for advanced profiles, while programmatic generation better preserves conservative safety and recovery constraints. The moderation should therefore not be read as a general LLM-source advantage.

The order result separates two quantities that are often conflated in LLM-as-judge work. One can estimate an order-balanced aggregate preference by swapping AB/BA positions, but that does not imply that the judge is making a stable pairwise decision. In this study, many pair--judge--run units are position-consistent rather than source-consistent. The marker-by-order results extend this point: decomposition alone does not remove position sensitivity. Some criteria, such as clarity/actionability, plan coherence, and recovery safety, are dominated by Plan-A effects; others, such as quality session design, are comparatively stable and consistently programmatic-favoring; and training specificity reverses direction across order. The practical lesson is that order swaps should not be treated as a cosmetic robustness check, and marker-level analyses should report marker-by-order and marker-by-regime interactions rather than only aggregate marker win rates.

The study has important limitations. It measures LLM judge preference, not objective coaching quality, and lacks human expert ratings. It uses only Qwen and Gemma judges in one structured generation domain. Structural matching controls observable weekly structure but not all quality-relevant coaching logic or presentation differences. The leakage-filtered analysis is post hoc, and the marker categories are diagnostic labels rather than validated constructs; the marker-by-order reversals show that some marker labels do not correspond to stable judge behavior under order swaps. Judge temperature was fixed at \(T=0.0\), improving reproducibility but limiting scope. Finally, row-level confidence intervals in the appendix should not be read as independent-sample inference because the same matched pairs are judged repeatedly. These caveats limit claims about training-plan quality, but they strengthen the central evaluation claim: source-preference estimates require order swaps, repeated-judgment diagnostics, criterion-by-order decomposition, athlete-band stratification, uncertainty reporting, and checks for presentation leakage.

\section{Conclusion}

LLM judges in this setting do not prefer directly LLM-source plans over structurally matched programmatic controls: the LLM-source win rate is 36.53\% overall. More importantly, the 42.02 percentage-point Plan A swing is larger than the aggregate source effect, and marker-by-order and marker-by-band results show that criterion-level judgments can be both position-sensitive and regime-dependent. Source-preference studies should therefore require order swaps, repeated-judgment diagnostics, criterion-by-order decomposition, athlete-band stratification, uncertainty reporting, and leakage audits.

\clearpage
\bibliography{references}

\clearpage
\onecolumn
\appendix

% Prompt, marker, schema, and structural-score appendix.
\section{Pairwise Judge Prompt Interface}
\label{app:pairwise-prompt}
The primary pairwise evaluation used \texttt{trailtraining.llm.soft\_eval.compare\_plans}. Unlike the marker task, this interface did not expose a project-local literal prompt string. The archived reproducibility object therefore consists of: (i) the exact canonical masked JSON payload passed to the evaluator, (ii) the judge model configuration, and (iii) the returned fields \texttt{preferred}, \texttt{reasoning}, \texttt{plan\_a\_advantages}, and \texttt{plan\_b\_advantages}. The evaluator compared two source-masked plans labeled only as Plan A and Plan B. Source family, source model, plan IDs, fixture metadata, self-family indicators, and AB/BA order were not included in the judge-facing payload. The call used \texttt{SoftEvalConfig(enabled=True, reasoning\_effort=``none'', skip\_synthesis=True, parallel\_batches=False, temperature=0.0)}. Appendix file \path{appendix_artifacts/pairwise_prompt_interface.txt} records this interface, and \path{appendix_artifacts/sample_pairwise_judge_input.json} contains a complete archived pairwise payload.

\section{Marker-Level Judge Prompt}
\label{app:marker-prompt}
The marker-level evaluation used the following system prompt exactly.

\begin{lstlisting}
You are an expert running coach and careful evaluator.

You will compare two anonymized training plans, Plan A and Plan B.
The plans have been source-masked. Do not infer or mention whether either plan was generated by a model, a program, or a human.

Evaluate each marker independently. For every marker, choose:
- "plan_a" if Plan A is better on that marker,
- "plan_b" if Plan B is better on that marker,
- "tie" if they are meaningfully equivalent.

Use only the visible plan content. Do not use source labels, file names, IDs, or provenance.
Return valid JSON only. No markdown.
\end{lstlisting}

The user prompt template was:

\begin{lstlisting}
Compare Plan A and Plan B on the following markers.

MARKERS:
- plan_coherence: Overall training-plan coherence: whether the week fits together as a coherent training plan, with compatible sessions, sensible sequencing, and no internal contradictions.
- training_specificity: Specificity to the athlete and trail-running context: terrain, hills, long-run needs, race phase, and appropriate trail/endurance demands.
- load_progression: Appropriateness of training load and progression: total volume, session duration, stress distribution, and whether the plan avoids abrupt or excessive load.
- recovery_safety: Recovery, rest, and safety: rest-day placement, hard-day spacing, injury-risk control, and avoidance of unsafe combinations.
- quality_session_design: Design of quality sessions: workouts such as intervals, tempo, hills, progression runs, or long runs are purposeful, specific, and feasible.
- endurance_development: Support for aerobic/endurance development: easy volume, long-run stimulus, sustainable aerobic work, and endurance-oriented structure.
- readiness_alignment: Alignment with readiness/recovery constraints: whether the plan respects the athlete's readiness, recovery capacity, and fatigue constraints.
- clarity_actionability: Clarity and actionability: whether the athlete can understand and execute the plan without ambiguity.
- explanation_quality: Quality of explanation/rationale: whether the stated purpose and explanation are helpful, specific, and support the plan. This marker is presentation-sensitive.

Scoring guidance:
- Scores are 1 to 5, where 5 is best.
- The preferred plan should normally be the one with the higher score.
- Use tie when differences are negligible.
- Keep rationales short.

REQUIRED JSON SCHEMA:
{"markers":{"plan_coherence":{"preferred":"plan_a | plan_b | tie","plan_a_score":"integer 1-5","plan_b_score":"integer 1-5","confidence":"number 0.0-1.0","rationale":"brief reason, max 18 words"}, ...},"overall_marker_summary":"one short sentence"}

PLAN INPUT:
{compact_json(judge_input)}
\end{lstlisting}
A complete instantiated marker prompt for one archived comparison is included as \path{appendix_artifacts/sample_complete_marker_prompt.txt}.

\section{Canonical Judge-Facing Schema}
\label{app:schema}
Both sources are rendered into the same canonical judge-facing schema. The judge sees only neutral \texttt{plan\_a} and \texttt{plan\_b} objects.

\begin{lstlisting}
{
  "plan_a": {
    "judge_view": "canonical_masked_structural_v1",
    "plan": {
      "days": [
        {
          "day_index": "integer 1..7",
          "session_type": "string",
          "duration_minutes": "integer",
          "is_rest_day": "boolean",
          "is_hard_day": "boolean",
          "target_intensity": "string",
          "workout_summary": "short normalized text",
          "purpose_summary": "short normalized text"
        }
      ],
      "summary": {
        "plan_days": "integer",
        "total_minutes": "integer",
        "rest_days": "integer",
        "hard_days": "integer",
        "long_runs": "integer"
      }
    }
  },
  "plan_b": "same schema as plan_a"
}
\end{lstlisting}
The source family, judge family, source model, plan IDs, structural-score gap, self-family flag, AB/BA order, judge identity, and run index are retained in the manifest outside the prompt for analysis. The frozen artifact set is named \texttt{canonical\_masked\_scrubbed\_v1\_t000}; the judge-facing payload inside each archived JSON record uses the \texttt{canonical\_masked\_structural\_v1} view shown above.

\section{Structural-Score Feature Definitions}
\label{app:structural-features}
The structural score excludes title wording, prose richness, citations, claim attributions, data-note verbosity, rationale or explanation fields, source model, generation arm, and file names. It uses only plan structure and fixture-conditioned expectations. The scored features are:

\begin{itemize}
  \item \textbf{plan days}: actual number of plan days compared with expected length;
  \item \textbf{total minutes}: sum of daily \texttt{duration\_minutes};
  \item \textbf{active, rest, hard, quality days}: counts derived from rest and hard flags and quality session types;
  \item \textbf{long runs}: count of \texttt{session\_type == long};
  \item \textbf{maximum and mean day minutes}: daily load distribution summaries;
  \item \textbf{hard-day spacing}: minimum distance between hard days;
  \item \textbf{fixture-adjusted load, hard-day, and rest-day expectations}: ranges conditioned on athlete band, readiness, recovery capability, and race phase.
\end{itemize}

The matching distance uses weighted gaps in total minutes, number of rest days, number of hard days, number of active days, number of long runs, number of quality days, maximum day minutes, and mean day minutes. Soft calipers penalize large gaps in total minutes, maximum day minutes, rest days, hard days, long runs, and quality days. Hard calipers were not enabled in the frozen primary match because the target-cardinality design prioritized retaining 250 same-cell pairs. The machine-readable definition is included as \path{appendix_artifacts/structural_score_features.json}.

\section{Sample Judge-Facing JSON}
\label{app:sample-json}
The following excerpt shows one canonical comparison. The full JSON appears in \path{appendix_artifacts/sample_pairwise_judge_input.json}.

\begin{lstlisting}
{
  "plan_a": {
    "judge_view": "canonical_masked_structural_v1",
    "plan": {
      "days": [
        {
          "day_index": 1,
          "duration_minutes": 48,
          "is_hard_day": false,
          "is_rest_day": false,
          "purpose_summary": "Build aerobic consistency while keeping fatigue low.",
          "session_type": "easy",
          "target_intensity": "easy",
          "workout_summary": "48 min easy run at conversational effort."
        },
        {
          "day_index": 2,
          "duration_minutes": 39,
          "is_hard_day": false,
          "is_rest_day": false,
          "purpose_summary": "Build aerobic consistency while keeping fatigue low.",
          "session_type": "easy",
          "target_intensity": "easy",
          "workout_summary": "39 min easy run at conversational effort."
        }
      ],
      "summary": {
        "hard_days": 0,
        "long_runs": 0,
        "plan_days": 7,
        "rest_days": 2,
        "total_minutes": 229
      }
    }
  },
  "plan_b": {
    "judge_view": "canonical_masked_structural_v1",
    "plan": {
      "days": [
        {
          "day_index": 1,
          "duration_minutes": 38,
          "is_hard_day": false,
          "is_rest_day": false,
          "purpose_summary": "Build confidence and endurance.",
          "session_type": "easy",
          "target_intensity": "easy",
          "workout_summary": "Easy trail run focusing on form and cadence."
        },
        {
          "day_index": 2,
          "duration_minutes": 41,
          "is_hard_day": false,
          "is_rest_day": false,
          "purpose_summary": "Maintain aerobic base and hill strength.",
          "session_type": "easy",
          "target_intensity": "easy",
          "workout_summary": "Easy trail run to maintain aerobic base."
        }
      ],
      "summary": {
        "hard_days": 0,
        "long_runs": 0,
        "plan_days": 7,
        "rest_days": 1,
        "total_minutes": 321
      }
    }
  }
}
\end{lstlisting}

\section{Sample Metadata Retained Outside the Prompt}
The manifest retains analysis metadata that the judge does not see. One example is shown below; the full file appears as \path{appendix_artifacts/sample_pair_metadata_outside_prompt.json}.

\begin{lstlisting}
{
  "record_id": "pair_0001__qwen_14b_judge__r00__BA",
  "pair_id": "pair_0001",
  "fixture_id": "ab_A1__r_high__rc_high__ph_base",
  "judge": "qwen_14b_judge",
  "judge_model_family": "qwen",
  "run": 0,
  "order": "BA",
  "left_plan_role": "programmatic",
  "right_plan_role": "llm",
  "llm_plan_id": "ab_A1__r_high__rc_high__ph_base__qwen2.5_7b_instruct__src_t070__exp_t000__s004",
  "programmatic_plan_id": "ab_A1__r_high__rc_high__ph_base__prog__exp_t000__s008",
  "source_model_family": "qwen",
  "self_family_match": true,
  "structural_score_gap": 0.0,
  "structural_score_version": "structural_score_v1.0.0",
  "feature_gaps": {
    "max_day_minutes": 35.0,
    "mean_day_minutes": 13.143,
    "n_active_days": 1.0,
    "n_hard_days": 0.0,
    "n_long_runs": 0.0,
    "n_quality_days": 0.0,
    "n_rest_days": 1.0,
    "total_minutes": 92.0
  },
  "judge_view": "canonical_masked"
}
\end{lstlisting}




\section{Sample Judge-Facing JSON and Metadata}
\label{app:samples}
Full examples are archived in \path{appendix_artifacts/sample_pairwise_judge_input.json} and \path{appendix_artifacts/sample_pair_metadata_outside_prompt.json}. The main report omits the full JSON listing to keep the page-limited portion focused on the experimental design and results.

\clearpage
\section{Detailed Results Tables}
\label{app:detailed-tables}
\renewcommand{\thetable}{\thesection.\arabic{table}}
\setcounter{table}{0}

\begin{table}[H]
\centering
\scriptsize
\setlength{\tabcolsep}{3pt}
\begin{tabular}{p{0.39\columnwidth}p{0.53\columnwidth}}
\toprule
Finding & Evidence \\
\midrule
No aggregate LLM-source preference & LLM-source wins 36.53\%; programmatic wins 63.47\%. \\
Athlete band moderates the effect & A1: 12.50\%; A4: 53.32\%. \\
Order dominates interpretation & 57.54\% as Plan A vs. 15.52\% as Plan B. \\
Markers are order/regime sensitive & Specificity is 33.50\% as A vs. 63.33\% as B; A4 specificity is 59.09\%. \\
Leakage filtering is stable & Clean subset LLM win rate is 34.44\%. \\
\bottomrule
\end{tabular}
\caption{Compact summary of the main results. Detailed results are reported in the remaining Appendix tables.}
\label{tab:summary}
\end{table}

\begin{table}[H]
\centering
\scriptsize
\setlength{\tabcolsep}{3pt}
\begin{tabular}{lrrr}
\toprule
Judge & LLM as A & LLM as B & Swing \\
\midrule
Gemma 12B & 39.60\% & 19.60\% & 20.00 pp \\
Gemma 4B & 62.00\% & 4.88\% & 57.12 pp \\
Qwen 14B & 74.40\% & 8.40\% & 66.00 pp \\
Qwen 7B & 54.16\% & 29.20\% & 24.96 pp \\
Overall & 57.54\% & 15.52\% & 42.02 pp \\
\bottomrule
\end{tabular}
\caption{Order effect by judge. Rates are LLM-source win rates when the LLM-source plan appears as Plan A versus Plan B.}
\label{tab:order-main}
\end{table}

\begin{table}[H]
\centering
\scriptsize
\setlength{\tabcolsep}{3pt}
\begin{tabular}{p{0.62\columnwidth}r}
\toprule
Diagnostic & Rate \\
\midrule
Source-consistent programmatic winner & 39.46\% \\
Source-consistent LLM-source winner & 12.52\% \\
Position-consistent Plan A winner & 45.02\% \\
Training specificity, LLM win rate excl. ties & 48.54\% \\
Load/quality/endurance, LLM win rates excl. ties & 35.11\% / 33.44\% / 32.81\% \\
Explanation quality, LLM win rate excl. ties & 32.21\% \\
\bottomrule
\end{tabular}
\caption{Main stability and marker diagnostics. Full pair-run consistency and marker tables appear in Appendix Tables~\ref{tab:pair-run-app},~\ref{tab:markers-app},~\ref{tab:marker-order-app}, and~\ref{tab:marker-band-app}.}
\label{tab:diagnostics-main}
\end{table}


\begin{table}[H]
\centering
\small
\begin{tabular}{lrr}
\toprule
Winner & Count & Rate \\
\midrule
LLM-source & 3,653 & 36.53\% \\
Programmatic & 6,347 & 63.47\% \\
\bottomrule
\end{tabular}
\caption{Primary forced pairwise outcome.}
\label{tab:primary-app}
\end{table}

\begin{table}[H]
\centering
\small
\begin{tabular}{llrrr}
\toprule
Band & Profile & Pairs & Primary & Marker \\
\midrule
A1 & Beginner & 60 & 12.50\% & 27.16\% \\
A2 & Developing & 55 & 24.55\% & 28.89\% \\
A3 & Advanced & 74 & 51.08\% & 44.45\% \\
A4 & Competitive & 61 & 53.32\% & 45.96\% \\
\bottomrule
\end{tabular}
\caption{LLM-source win rates by athlete band. Primary is the forced pairwise win rate. Marker is the LLM-source win rate excluding ties across all marker decisions.}
\label{tab:athlete-app}
\end{table}

\begin{table}[H]
\centering
\small
\begin{adjustbox}{max width=\textwidth}
\begin{tabular}{lrrrr}
\toprule
Marker & A1 & A2 & A3 & A4 \\
\midrule
Clarity/actionability & 29.79\% & 31.88\% & 46.95\% & 50.95\% \\
Plan coherence & 25.23\% & 27.00\% & 43.38\% & 45.21\% \\
Recovery safety & 40.07\% & 40.53\% & 40.23\% & 35.60\% \\
Readiness alignment & 41.09\% & 42.89\% & 39.09\% & 37.46\% \\
Load progression & 26.82\% & 28.78\% & 40.44\% & 42.35\% \\
Endurance development & 16.88\% & 17.55\% & 44.19\% & 49.38\% \\
Quality session design & 11.50\% & 17.86\% & 46.92\% & 51.06\% \\
Explanation quality & 23.16\% & 19.92\% & 39.05\% & 42.11\% \\
Training specificity & 35.58\% & 37.02\% & 58.72\% & 59.09\% \\
\bottomrule
\end{tabular}
\end{adjustbox}
\caption{Marker-level LLM-source win rates across athlete bands, excluding ties. Each cell reports the percentage of non-tied marker decisions won by the LLM-source plan. A1 is the beginner/low-volume band and A4 is the competitive/high-volume band.}
\label{tab:marker-band-app}
\end{table}


\begin{table}[H]
\centering
\small
\begin{tabular}{lrrr}
\toprule
Judge & LLM as Plan A & LLM as Plan B & Plan-A advantage \\
\midrule
Gemma 12B & 39.60\% & 19.60\% & 20.00 pp \\
Gemma 4B & 62.00\% & 4.88\% & 57.12 pp \\
Qwen 14B & 74.40\% & 8.40\% & 66.00 pp \\
Qwen 7B & 54.16\% & 29.20\% & 24.96 pp \\
Overall & 57.54\% & 15.52\% & 42.02 pp \\
\bottomrule
\end{tabular}
\caption{Order effect by judge. Rates are LLM-source win rates when the LLM-source plan appears as Plan A versus Plan B.}
\label{tab:order-effect-app}
\end{table}

\begin{table}[H]
\centering
\small
\begin{tabular}{lrr}
\toprule
Pair-run category & Count & Rate \\
\midrule
Source-consistent programmatic & 1,973 & 39.46\% \\
Source-consistent LLM-source & 626 & 12.52\% \\
Position-consistent Plan A & 2,251 & 45.02\% \\
Position-consistent Plan B & 150 & 3.00\% \\
\bottomrule
\end{tabular}
\caption{AB/BA pair-run consistency over 5,000 matched pair--judge--run units.}
\label{tab:pair-run-app}
\end{table}

\begin{table}[H]
\centering
\small
\begin{tabular}{lrrrr}
\toprule
Marker & LLM wins & Prog. wins & Ties & LLM win rate excl. ties \\
\midrule
Training specificity & 4,368 & 4,630 & 1,002 & 48.54\% \\
Clarity/actionability & 2,710 & 3,962 & 3,328 & 40.62\% \\
Readiness alignment & 2,863 & 4,324 & 2,813 & 39.84\% \\
Recovery safety & 3,103 & 4,839 & 2,058 & 39.07\% \\
Plan coherence & 2,835 & 5,040 & 2,125 & 36.00\% \\
Load progression & 3,218 & 5,947 & 835 & 35.11\% \\
Quality session design & 2,783 & 5,539 & 1,678 & 33.44\% \\
Endurance development & 3,046 & 6,239 & 715 & 32.81\% \\
Explanation quality & 2,578 & 5,425 & 1,997 & 32.21\% \\
\bottomrule
\end{tabular}
\caption{Marker-level outcomes across 10,000 records per marker. The final column excludes ties; tie counts are shown because the marker task permits equivalence judgments.}
\label{tab:markers-app}
\end{table}

\begin{table}[H]
\centering
\small
\begin{adjustbox}{max width=\textwidth}
\begin{tabular}{lrrr}
\toprule
Marker & LLM as Plan A & LLM as Plan B & Plan-A minus Plan-B \\
\midrule
Clarity/actionability & 85.40\% & 1.21\% & 84.19 pp \\
Plan coherence & 76.16\% & 1.30\% & 74.86 pp \\
Recovery safety & 76.03\% & 5.92\% & 70.11 pp \\
Readiness alignment & 59.55\% & 22.86\% & 36.69 pp \\
Load progression & 46.63\% & 23.44\% & 23.19 pp \\
Endurance development & 39.82\% & 25.48\% & 14.34 pp \\
Quality session design & 34.04\% & 32.85\% & 1.19 pp \\
Explanation quality & 24.02\% & 39.76\% & -15.74 pp \\
Training specificity & 33.50\% & 63.33\% & -29.83 pp \\
\bottomrule
\end{tabular}
\end{adjustbox}
\caption{Marker-level LLM-source win rates by order, excluding ties. Positive values in the final column mean the marker is more favorable to LLM-source plans when the LLM plan is shown as Plan A; negative values mean it is more favorable when the LLM plan is shown as Plan B.}
\label{tab:marker-order-app}
\end{table}


\begin{table}[H]
\centering
\small
\begin{tabular}{lrrrrr}
\toprule
Comparison & $n$ & LLM wins & Rate & CI low & CI high \\
\midrule
Overall & 10,000 & 3,653 & 36.53\% & 35.59\% & 37.47\% \\
LLM as Plan A & 5,000 & 2,877 & 57.54\% & 56.17\% & 58.91\% \\
LLM as Plan B & 5,000 & 776 & 15.52\% & 14.52\% & 16.52\% \\
\bottomrule
\end{tabular}
\caption{Descriptive row-level Wald 95\% confidence intervals. These are not independent-sample inference because rows repeat matched pairs across judges, runs, and orders.}
\label{tab:uncertainty-app}
\end{table}

\begin{table}[H]
\centering
\small
\begin{tabular}{lrrrrr}
\toprule
Sample & Pairs & Judgments & LLM wins & Prog. wins & LLM win rate \\
\midrule
Full sample & 250 & 10,000 & 3,653 & 6,347 & 36.53\% \\
Artifact-clean subset & 103 & 4,120 & 1,419 & 2,701 & 34.44\% \\
Artifact-present subset & 147 & 5,880 & 2,234 & 3,646 & 37.99\% \\
\bottomrule
\end{tabular}
\caption{Full-sample and leakage-filtered sensitivity analysis. The artifact-clean subset excludes any matched pair flagged by the post hoc presentation-artifact audit.}
\label{tab:leakage-app}
\end{table}


\end{document}
